\documentclass[11pt,a4paper]{article}
\usepackage[T1]{fontenc}
\usepackage{lmodern}
\usepackage[margin=25mm]{geometry}
\usepackage{amsmath,amssymb,amsthm,mathtools}
\usepackage{microtype}
\usepackage{booktabs,tabularx,array}
\usepackage{xcolor}
\usepackage[colorlinks=true,linkcolor=blue!40!black,citecolor=blue!40!black,
  urlcolor=blue!40!black]{hyperref}
\hypersetup{pdftitle={A Single-Exponential Exponent for Arithmetic Progressions},
  pdfsubject={Research draft},
  }
\setlength{\parindent}{1.2em}
\setlength{\parskip}{2pt}
\setlength{\emergencystretch}{2em}
\allowdisplaybreaks[1]
\newtheorem{theorem}{Theorem}[section]
\newtheorem{lemma}[theorem]{Lemma}
\newtheorem{proposition}[theorem]{Proposition}
\theoremstyle{definition}
\newtheorem{definition}[theorem]{Definition}
\newcommand{\M}{\mathcal M}
\newcommand{\E}{\mathbb E}
\newcommand{\Z}{\mathbb Z}
\newcommand{\R}{\mathbb R}
\newcommand{\g}{\mathfrak g}
\newcommand{\h}{\mathfrak h}
\newcommand{\n}{\mathfrak n}
\newcommand{\Ad}{\operatorname{Ad}}
\newcommand{\ev}{\operatorname{ev}}
\newcommand{\Lip}{\operatorname{Lip}}
\newcommand{\wt}{\operatorname{wt}}
\newcommand{\corr}{\operatorname{corr}}
\newcommand{\poly}{\operatorname{poly}}
\newcommand{\im}{\operatorname{image}}
\newcommand{\one}{\mathbf 1}

\begin{document}
\begin{center}
{\Large\bfseries A Single-Exponential Exponent\\[3pt]
for Arithmetic Progressions}\\[7pt]
{\normalsize AI-assisted research draft\textsuperscript{*}}\\[5pt]
{\small 8 October 2026}
\end{center}
\begingroup
\renewcommand{\thefootnote}{\fnsymbol{footnote}}
\footnotetext[1]{AI assistance: OpenAI ChatGPT was used for mathematical exploration and drafting; OpenAI Codex was used for editing and compilation.}
\endgroup

\begin{abstract}
Using the analytic inputs of \emph{Quasipolynomial Bounds for Arithmetic
Progressions} (OpenAI/math, family 159), we obtain an admissible density
exponent \(\varepsilon_k=\exp(-k^C)\) for an absolute constant \(C\).
Two changes give the single-exponential dependence. An ordered return
factorization keeps historical substitutions out of the next active orbit,
so the geometric degree is \(O(k^2)\). An inverse-theorem induction retains
the supplied derivative witnesses, so its logarithmic-budget exponent is
single exponential in the order. The original rank bookkeeping and a
descending width recurrence then suffice. Only the single-exponential
dependence on \(k\) is tracked.
\end{abstract}

\section{The result and its inputs}

Let \(r_k(N)\) be the largest size of a subset of \([N]\) containing no
nonconstant \(k\)-term arithmetic progression. Write \(\M\) for the
September 23, 2026 manuscript \cite{openai}.

\begin{theorem}\label{thm:main}
Assume the analytic inputs of \(\M\) listed below. There is an absolute
constant \(C>0\) such that, for every \(k\ge3\), there are constants
\(K_k,c_k>0\) for which
\[
 r_k(N)\le K_kN\exp\!\left[-c_k(\log N)^{\varepsilon_k}\right],
 \qquad
 \varepsilon_k=\exp(-k^C),\qquad N\ge3.
\]
\end{theorem}

The inputs are the bottom absolute-patch theorem (Lemma 2.2), the
sampler and rank-transfer estimates (Sections 3--4 and Appendices F--H),
the field and active-return framework (Definition 8.1 and Theorem 8.2),
the exact macro extraction (Section 9), and the compression and
descending numerical schedule (Section 10). The inverse-theorem input is
the argument of Leng--Sah--Sawhney \cite{lss}, with its quantitative
induction modified in Section~\ref{sec:numerical}. We prove the additional
algebraic and quantitative assertions required for these modifications.
The conclusion retains this dependence on the analytic inputs of \(\M\).

Throughout, constants depending only on \(k\), or on a fixed theorem
order, may be absorbed into multiplicative constants. Powers of growing
budgets are retained. Thus a bound \(K_Dx^{\exp(D^{O(1)})}\) is acceptable,
whereas an uncontrolled iteration of such a power is not. Every
\(O(1)\) in an exponent below is absolute.

For a starting density \(\alpha>0\), put
\[
 s=k-2,\qquad D=256k^2,\qquad p=2+\log(1/\alpha).
\]
The density increment takes \(T\ll_k p\) rounds. At weight \(h\), the
determining data have integer coordinates \(b_h\), a rational value space
\(W_h\), and targets
\[
 C_h(u,b_{<h})\in W_h,\qquad \deg_{\wt} C_h\le h.
\]
The spatial variable has weight one and \(b_h\) has weight \(h\).
The strict and enlarged determining cutoffs force a unique integer tuple.
If \(w_h\) is the width, write \(Q_h=\log(2/w_h)\). The essential order
condition is that a width loss at \(h\) may depend on numerical data
already fixed at weights \(>h\), but not on incoming \(Q_i\) with
\(i\le h\). We call such a loss \emph{warm at \(h\)}. Data using the
current \(Q_h\) are allowed in a preliminary bound at \(h\); when
retained as history, they are warm at every later weight \(j<h\).

The proof has three quantitative outputs:
\[
 \text{geometric degree }D=O(k^2),\qquad
 \text{local budget degree }B\le\exp(D^{O(1)}),\qquad
 \varepsilon_k^{-1}\le\exp(D^{O(1)}).
\]
We keep the original rank-cut bookkeeping; no improvement to that
bookkeeping is needed.

\section{A persistent record with quadratic degree}\label{sec:macro}

A completed active return has two kinds of controlled factors: a slow
factor in normalized coordinates and a rational factor in integer
coordinates. Keeping these factors on their original sides allows the
middle orbit to retain its original degree.

For a filtered Lie group, call a polynomial \emph{jointly filtered} when
each monomial of weighted degree \(a\) in its logarithm has coefficient
in the \(a\)-th Lie layer. Constants lie in the first layer. Integer
and normalized labels have weight one. A \emph{marking} is the fixed
projection \(\pi:G\to Y\) recording the determining equations.

\begin{lemma}[Ordered active return]\label{lem:active}
In the setting of Theorem 8.2 of \(\M\), suppose the input orbit is
jointly filtered in all actual parameters and its marking is a fixed
filtered polynomial \(P_Y(X)\) of the site variables. The group and
marking may themselves be finite charged families in older dummy
variables, which are coefficients rather than actual arguments.
The return can be chosen in the form below, with each of \(E,C,R\)
jointly filtered in the original filtration:
\begin{equation}\label{eq:ECR}
 g^\sharp(X,\theta,\widetilde t,z)
   =E(X,\theta,z)\,C(X,z)\,R(X,\widetilde t,z),
 \qquad
 \pi E=\pi R=1,\quad \pi C=P_Y(X).
\end{equation}
Here \(X\) collects the site variables, \(z\) the short labels, and
\(\theta,\widetilde t\) the normalized and integer long labels.
The factor \(C\) has no long-label argument; \(E\) has controlled coefficients
in normalized site, residual, and label coordinates; and \(R\) has
controlled common denominator in raw integer coordinates. No
coefficient bound is required for \(C\). The original score slice,
both expected-payoff representative identities, and the field
estimates are retained.
\end{lemma}

The additional induction proving the lemma is given in
Section~\ref{sec:active}. We first show how its format closes the macro
argument.

\begin{proposition}[Persistent macro record]\label{prop:macro}
Starting from an absolute patch of degree \(s\) and rank \(d_0\), the
macro argument can retain the original \(d_0\) absolute slots and a
middle orbit of joint degree at most \(s\). Its final physical formula
has absolute-weight ceiling at most \(s(s+1)\). The representation,
field, score, and extraction contracts of \(\M\) are preserved, with
warm costs at each later layer. The padded ceiling \(D=256k^2\)
therefore suffices.
\end{proposition}

\begin{proof}
\textit{The retained data.}
Store the return as \(E_{\rm hist}CR_{\rm hist}\), and use only \(C\)
as the next active orbit. Historical integer short labels remain
formal coefficient variables, with \emph{one total-charge cap} \(s\)
over all their blocks. Raw and normalized copies of a short label
always use the same variable and its fixed normalizing divisor.
Historical long labels remain in the outer factors. A deep
coefficient lattice evaluates into the pointwise lattice at every
integer short label; arbitrary real labels have no such property.
Before the next positive-filtration recursion, normalize any
degree-zero family constant as in the proof of
Lemma~\ref{lem:active}.

A newly promoted determining coordinate retains its original weight
in \(C\). In outer-factor coefficients use its canonical site value:
the left factor uses the bounded normalized residual frame, and the
right factor uses raw integer coordinates. Exact marking identifies
these values at every tested point. Fix the integral determining
reindexings before recording history, and keep their exact
triangular inverses. They preserve the original determining degrees.
Later congruence restrictions have their creation-layer costs, hence
warm costs at every lower layer.

\textit{Ordered representation.}
Normalize the middle factor only. For an integral triangular
reindexing \(T\) of the absolute tuple, the correct expression is
\begin{equation}\label{eq:reindex}
 E_{\rm hist}(CT)(T^{-1}R_{\rm hist}T).
\end{equation}
This is the old physical map followed by \(T\), which permutes the
complete integer tuple sum. Conjugation by the integral map preserves
controlled rational denominators; its numerator need not be bounded.
Coefficientwise reduction gives \(CT=Sg\), where \(S\) is the bounded
full-weight map. It is independent of the historical variables:
their positive degree excludes them from a full-weight slot term.
If \(g=E_{\rm new}C_{\rm new}R_{\rm new}\), collect the result as
\begin{equation}\label{eq:collect}
 \bigl(E_{\rm hist}\Ad_S E_{\rm new}\bigr)
 (SC_{\rm new})
 \bigl(R_{\rm new}R'_{\rm hist}\bigr),
 \qquad R'_{\rm hist}=T^{-1}R_{\rm hist}T.
\end{equation}
Only bounded maps conjugate slow factors. Rational right factors
stay on the right. The triangular subgroup in which determining
transformations have no absolute arguments is bracket closed;
forgetting the absolute slots is strong, since a determining shear
lifts while fixing them. Thus this normalization preserves the
exact determining projection.

\textit{The quotient fields.}
At site \(u\), define the field on the coefficient-family quotient by
the original cutoff sum evaluated on
\begin{equation}\label{eq:field}
 E_{\rm site}(u)\,S\,
       \ev_{w(u)}(Y)\,R_{\rm site}(u).
\end{equation}
Canonical site substitution occurs only in the external coefficients
and in \(\ev_{w(u)}\). All original determining and absolute integer
variables are still summed. On an exactly marked orbit, any nonzero
determining cutoff forces the unique canonical tuple. Hence both
payoffs agree exactly with the originals, including their zero
cases.

A common deep coefficient lattice, further deepened for the rational
right factors, makes \eqref{eq:field} well defined. On bounded group
representatives, the left map and \(S\) are bounded; the right map
can be reduced to a bounded right-coset representative determined by
its denominator and coefficient residues. The triangular inverse
in the representation argument then bounds all contributing integer
tuples and derivatives. Bounded-range integer evaluation has
controlled Lipschitz norm. This gives the required Lipschitz bound,
independently of the unrestricted coefficients of \(C\).
Triangular uniqueness preserves the cap one.

\textit{The two field meshes.}
After extraction, product-column labels are frozen. The only
historical short labels that remain variable are integral affine
maps from extra inactive rows. Those rows annihilate the prepared
value space. The chart estimate has the form
\[
 \|z_i\|\le \exp(\poly(d))K_i,\qquad K_i\le k_i/v,
\]
where \(v\) is the chart width. For a short axis, its scale \(k_i\),
the integral offset, and its values throughout the fixed padded
chart have preliminary bounds. The factor \(v^{-1}\) is included
in that bound. At every later layer all these costs are warm.
The fixed determining polynomials therefore turn the finite
row-and-offset menu into a finite menu of entire functions
\(u\mapsto w(u)\).

Historical real labels are affine expressions in the current
normalized spatial coordinates and the fixed higher normalized
residuals. They are bounded on the same padded chart, even when
they leave a formerly selected unit interval. Retain the original
smooth cutoffs and score slice; no new moving sharp indicator is
inserted. Outside this fixed padded chart use the zero extension
from the representation construction of \(\M\).

At one site, mesh bounded left values, full-weight coefficients,
bounded integer evaluations, right cosets, and inherited cutoff
parameters. For a mesh uniform over all sites, mesh the bounded
left \emph{coefficient vectors}, the finite whole-function short
menu, rational coefficient residues, and the inherited field
parameters. The canonical site functions are fixed throughout.
These give the distinct point-value and full-field meshes required
by Definition 8.1 of \(\M\).

First pigeonhole the finite tuple of actual historical denominators
and cover moduli, then take their product and a fixed deepening.
Its logarithm is a sum of warm budgets. Taking the least common
multiple of all integers below an upper bound is unnecessary.
After product-column freezing, historical raw labels are affine
in raw spatial and determining integers and the remaining extra
inactive labels. Later integer sections may have arbitrarily large coefficients;
they change only the rational numerators, which are recorded
modulo the common deep grid. The same divisibility gives the
long-label periods. Selection pays a warm loss in the fraction of
available productive paths; every retained path keeps its original
full score and slice.

\textit{Extraction and the final degree.}
For a macro comparison, reassemble the physical formula. On its
scored locus, replace canonical integers in external coefficients
by their corresponding integer summation variables. Exact marking
and the original cutoffs identify the expressions. Current
auxiliary labels now enter only through this polynomial formula,
and both expected representative identities apply before the
stochastic constants are frozen. Section 9 of \(\M\) uses these
identities, the original slice and masks, cap one, exact
determining and inactive equations, and a degree ceiling. It
does not require bounds on coefficients or spatial derivatives
of the returned physical polynomial. Its extraction and
localization therefore apply with the original payoffs.

Every recorded historical-label substitution has final
spatial/determining weighted degree at most \(s\). For a monomial
of original absolute-input weight \(h'\) and other charge \(d\),
the joint invariant gives \(h'+d\le h\), where \(h\le s\) is
its absolute destination weight. Give each absolute slot the
temporary weight \((s+1)h\). Simultaneous substitution gives
\begin{equation}\label{eq:quadratic}
 (s+1)h'+sd\le (s+1)(h'+d)\le(s+1)h.
\end{equation}
This reweights absolute inputs and destinations together.
Consequently the physical ceiling is \(s(s+1)\).

These temporary weights are used for extraction and final
compression. The next middle orbit still uses the original
weights, with determining slots and historical integer variables
formal; its outer factors again enter the fields. This closes
the induction without repeated degree inflation. No family or
partner coordinate becomes an absolute slot, so the root rank
is still \(d_0\), independent of old numerical widths.

Finally, \(m\) historical variables of total charge at most \(s\)
have at most
\[
 \binom{m+s}{s}\le(m+1)^s
\]
monomials per scalar coordinate. There are at most \(s\) macro
layers and \(s\) internal degree reductions per return.
Rational systems, coefficient arrays, and finite meshes are
handled simultaneously. The added operations have polynomial
sequential depth in \(s\); their quantitative bounds are included
in Section~\ref{sec:numerical}.
\end{proof}

\section{The charged active return}\label{sec:active}

We prove Lemma~\ref{lem:active} by supplementing the degree induction
of Theorem 8.2 of \(\M\). The point is to retain total degree when
short variables become coefficient variables, while each recursive
child still has strictly smaller positive degree.

\begin{proof}[Proof of Lemma~\ref{lem:active}]
\textit{Charged families and positive normalization.}
Record rational coefficient spaces \(A_{d,\ell}\) with
\[
 [A_{d,\ell},A_{e,m}]\subset A_{d+e,\ell+m},\qquad
 A_{0,\ell}=A_{1,\ell},\qquad A_{r+1,\ell}=0,
\]
and one finite total-charge cutoff in \(\ell\).
Initially \(A_{d,\ell}=\g_{\max(d,\ell)}\), with
\(\g_0=\g_1\). Promoting a short block to a dummy \(w\) replaces
the array by
\begin{equation}\label{eq:family}
 A'_{d,\ell}
   =\left\{\sum_\beta w^\beta a_\beta:
                 a_\beta\in A_{d,\ell+|\beta|}\right\}.
\end{equation}
The positive layers are \(\exp A'_{i,i}\), and the
degree-zero group is \(\exp A'_{0,0}\). Quotients use images
of these spaces, and refinements use the actual intersections
before taking images. Every raw and normalized occurrence of a
promoted label is changed simultaneously, including those in
marks and ancestor factors.

Perform the common-site marking correction before fast-algebra
selection. A coefficientwise ordinary filtered section commutes
with evaluation at the actual short label, so the inserted
correction evaluates to the identity. Normalize the fixed marked
constant by a bounded left factor and a rational right factor.
The normalized mark takes values in the positive marked group.
Let \(c=g(0,0;w)\) be the constant family element after this
marked normalization.
The remaining constant coset modulo the positive group can now
be removed in the kernel: choose a positive-layer lift \(c_1\)
of its marked image, and split the difference
\(n_0=cc_1^{-1}=e_N\gamma_N\) into a bounded kernel element
and a kernel lattice element. Replace the family orbit \(g\) by
\(e_N^{-1}g\gamma_N^{-1}\), retaining both factors for restoration.
This lies in the positive group because its constant coset
does. These are family-group operations, so all charge bounds
are retained. A deep source lattice with the prescribed image
handles the marked rational factor; its numerator need not be
bounded.

\textit{Refiltration with the old degree retained.}
The degree-zero action is made compatible with the refiltration
by adjoining to the detected fast algebra the rational graded
ideal
\[
 S_i=\im\!\left(A'_{i+1,i}\longrightarrow
                    A'_{i,i}/A'_{i+1,i+1}\right).
\]
The bracket rule makes degree-zero conjugation trivial modulo
\(S_i\), and \(S_r=0\). Thus saturation preserves the detected
top and its injective marked projection, while the resulting
refiltration is normalized by the degree-zero group.

Write \(G_i\) for the old positive layers and \(H_i\) for
the refiltration from Lemma 7.6 of \(\M\). The prescribed
full factors preserve old joint degree: a grade-\(i\) symbol
correction is lifted into \(G_i\); the off-space substitution
preserves weighted degree; and the constant adjustment sees
only old grade one, the higher short degrees already lying
in \(G_2\subset H_1\). A pure-long grade-\(i\) symbol is independent
of the short variables, since every coefficient of positive
short degree lies in old \(G_{i+1}\); its kernel lift therefore
preserves total degree. Filtered multiplication and conjugation
preserve these bounds. The refiltered input therefore satisfies
\begin{equation}\label{eq:mixed}
 [v^\alpha z^\beta]\log h
       \in \h_{|\alpha|}\cap\g_{|\alpha|+|\beta|},
 \qquad \h_0=\h_1.
\end{equation}
Here \(v\) denotes the long variables and \(z\) the current
short variables.

Since \(G_{i+1}\subset H_i\subset G_i\), the layers
\[
 H_1\supset G_2\supset H_2\supset G_3\supset\cdots
\]
form one rational flag, with its strong image in the marked
group. A rational section adapted to this flag preserves the
required intersections. It gives a fixed exact lift \(p(X,z)\)
of the refiltered mark with
\[
 [X^\gamma z^\beta]\log p\in
       \h_d\cap\g_{d+|\beta|},\qquad d=\wt\gamma.
\]
Only this interleaved local flag is used; historical charge is
already encoded in the coefficient-family group.

\textit{One lower-degree child.}
Let \(q:H\to Q=H/H_r\), and let \(N\) be the actual
kernel of the marking on \(H\). Top injectivity identifies
\(N\) with the kernel of the quotient marking via \(q|_N\).
Promote the current short block to a new dummy \(w\) and form
\begin{align*}
 \mathcal Q_i
   &=\left\{\sum_\beta w^\beta a_\beta:
           a_\beta\in q(\h_i\cap\g_{i+|\beta|})\right\},\\
 \mathcal N_i
   &=\left\{\sum_\beta w^\beta n_\beta:
           n_\beta\in q(\n\cap\h_i\cap\g_{i+|\beta|})\right\}.
\end{align*}
At index zero use \(\h_0=\h_1\) and \(\g_0=\g_1\).
These are images of the displayed intersections.
The bracket relation
\([\mathcal Q_i,\mathcal N_j]\subset\mathcal N_{i+j}\)
defines the split rational family
\(\mathcal N\rtimes\mathcal Q\), marked by its second
coordinate. Its canonical section is \(a\mapsto(1,a)\).
The map \((n,a)\mapsto(na,a)\) embeds it into two copies of the
quotient family, so it is a rational nilpotent group.
Both degree-\(r\) spaces vanish, so its positive group has
degree at most \(r-1\).

The local pair
\[
 \bigl(qh\,(qp)^{-1},\,qp\bigr)
\]
has the required kernel coefficients by \eqref{eq:mixed}.
Correct its marking to the common site with the canonical
section, and apply the degree-zero normalization above.
The resulting input is jointly filtered in all actual
variables: substituting a weight-\(d\) site monomial can
produce actual degree \(a\le d\), and a dummy degree \(b\)
is charged in the old layer \(\g_{d+b}\subset\g_{a+b}\).
Apply the induction to this single child, leaving \(w\)
unevaluated.

Restore the child's preliminary degree-zero factors, so its
exact second coordinate is \(qp(X,w)\). Multiplying by the
inverse of its canonical marked section isolates the kernel
family. Lift this family through \((q|_N)^{-1}\), multiply
by \(p(X,w)\), and evaluate \(w=z\). Its degree-\(a\),
dummy-degree-\(b\) coefficients lift into
\(\n\cap\h_a\cap\g_{a+b}\). Thus reconstruction restores
the original joint degree exactly, including synchronized
normalized copies. A common deep coefficient lattice
evaluates inside the pointwise lattice at every integer
\(z\). Fields pull back through the homomorphism
\((n,a)\mapsto na\); pullback preserves supremum meshes,
and bounded short evaluation has controlled Lipschitz cost.
No metric bound on the particular marked lift \(p\) is used.

\textit{Ordered factors and the original score.}
Keep the child return in the full split marked group.
The joint-fiber reconstruction \(T\) is equivariant for
actual kernel factors:
\[
 T\bigl((e,1)c(r,1)\bigr)=\widetilde e\,T(c)\,\widetilde r.
\]
Indeed the two sides have the same quotient and exact
mark, and the joint embedding is injective. The lifts
\(\widetilde e,\widetilde r\) are rational and filtered.
Collect them through the prescribed bounded left and
rational right factors of \(\M\). For the local notation
\(K_E\rho(E_F)\) and \(\rho(R_F)K_R\), this gives
\[
 E=K_E\Ad_{\rho(E_F)}(\widetilde e),\quad
 C=\rho(E_F)T(c)\rho(R_F),\quad
 R=\Ad_{\rho(R_F)^{-1}}(\widetilde r)K_R.
\]
The left conjugation has bounded normalized coefficients,
and the right one has controlled rational denominator.
The middle has no long-label argument.

If \(L\) is the marking restorer of \(\M\), use
\(E(LC)R\). It has the required marking everywhere and
agrees with \(LECR\) on the full exact-inactive forecasting
locus, where \(L=1\). That is the locus of both original
payoffs and their representative identities. At an earlier
central quotient, lift the outer factors separately in
the actual kernel and lift the middle separately.
Any difference from the original selected lift is removed
by the independent central Haar average. Restore the
initial degree-zero factors and project away partners.
The marking and joint degree are exact, and the expected
payoffs, original slice, and additive analytic errors
are those of \(\M\).

At degree zero the scalar-comparison base case applies.
Every positive step has one strictly smaller-degree child.
This proves the lemma. All additional operations are
simultaneous rational linear algebra, bounded-degree
coefficient constructions, and one split product per step.
\end{proof}

\section{Single-exponential numerical costs}\label{sec:numerical}

There are two separate numerical tasks. The external inverse
theorem must have a single-exponential \emph{budget exponent};
the operations inside \(\M\) must preserve that class.

\subsection{The external inverse exponent}

Let \(I_s\) be an exponent such that the degree-\(s\)
inverse theorem produces a correlating nilsequence with
\begin{equation}\label{eq:inverse-budget}
 d+\log M+\log(2+\|F\|_{\Lip})+\log(1/\corr)
 \le K_s\bigl(2+\log(1/\delta)\bigr)^{I_s}.
\end{equation}
Here \(d\) is the nilmanifold dimension, \(M\) its
complexity, and \(\delta\) the Gowers-norm threshold.

\begin{proposition}\label{prop:inverse}
The inverse-theorem argument can be run with
\[
 I_s\le\exp(s^{O(1)}).
\]
\end{proposition}

\begin{proof}
\textit{Dimension and logarithmic height.}
Keep \(d\) separate from the logarithmic numerical budget
\(\Lambda\), which includes rational heights, reciprocal
errors, Lipschitz norms, and the logarithm of the
nilcharacter's output dimension. Basic rational and
coordinate operations in order \(q\) have the form
\begin{equation}\label{eq:height}
 d_{\rm out}\le K_qd,\qquad
 \Lambda_{\rm out}\le K_q(2+d)^b\Lambda_{\rm in}.
\end{equation}
Common denominators are cleared jointly; maximal minors
give bases, kernels, and sections with logarithmic height
polynomial in \(d\) and linear in \(\Lambda\).
For bounded-step BCH, all bracket denominators divide a
fixed power of the common structure-constant denominator.
Fourier term counts and bounded-step products likewise
have logarithmic cost polynomial in \(d\) and linear in
\(\Lambda\).

In Leng's degree/step induction \cite{leng}, each recursive
group is a subgroup or a square-group quotient of
dimension at most \(2d\). The next recursive call uses the
new logarithmic budget but a dimension comparable to the
original \(d\). Thus its dimension exponents satisfy
\[
 b_{a,b}\le b_{a-1,b}+b_{a,b-1}+q^{O(1)},
 \qquad b_{a,1},b_{1,b}\le q^{O(1)}.
\]
Pascal induction gives \(b_{q,q}\le\exp(O(q))\).
The later step reductions of \cite[Theorem 5.4]{lss}
stay in subgroups or quotients and add such exponents.
The simultaneous denominator bounds in the rational
calculus handle the dimension loops without repeated
height substitution.

Take \(q=O(s^2)\) for the degrees and rank filtrations
appearing in \cite{lss}. Its universal groups have
dimension at most \(K_qm^q\) on \(m\) generators.
Polarization and the constructions for powers and roots
use an order-dependent number of copies, and the latter
use subgroups of \(G\times G\); they do not create a
dimension power equal to the integer root order.
Chart counts and nilcharacter output coordinates are
charged logarithmically. The finite helper chains,
degree-rank passages, and converse-inverse reductions
have polynomial sequential depth in \(s\). Consequently
the non-inverse part of an order-\(s\) construction has
budget degree
\begin{equation}\label{eq:Fs}
 F_s\le\exp(s^{O(1)})
\end{equation}
in \(2+d+\Lambda\).

The additive-combinatorial primitive and the fixed \(U^3\)
base case have polynomial logarithmic budgets: the
Bogolyubov--Ruzsa and progression-to-Bohr steps are
\cite{sanders,milicevic}, the quadratic construction is
\cite[Appendix A]{lengquadratic}, and its nilsequence
conversion is \cite{jamneshan}. In the last conversion,
the group has dimension \(O(m^2)\), the radius loss has
logarithm \(O(m\log m)\), and the defining cutoff sum
has at most \(3^m\) terms on a fundamental cube. Thus
its logarithmic Lipschitz cost is polynomial as well.
All these fixed-order costs can be included in \(F_s\).

\textit{Keeping derivative witnesses.}
We need a stronger induction than the ordinary inverse
statement. Let \(J_s\) assert that if a 1-bounded function
\(f:[N]\to\mathbb C\) has, for some \(H\subset[N]\)
with \(|H|\ge\rho N\) and \(0<\rho\le1\), supplied
degree-\((s-1)\) nilsequences
\(\nu_h\) such that
\[
 \left|\E_{n\in[N]}\Delta_h f(n)\nu_h(n)\right|\ge\rho,
 \qquad \Delta_hf(n)=f(n)\overline{f(n+h)},
\]
then \(f\) has a degree-\(s\) correlating nilsequence with
budget at most \(K_sx^{J_s}\), where
\[
 x=2+d+\log M+\log(1/\rho).
\]
We use the zero extension of \(f\). The witnesses may
have different nilmanifolds and functions, each of dimension
at most \(d\), complexity at most \(M\), and Lipschitz norm
at most \(M\). They are normalized to be 1-bounded.

The elementary step is the following consequence of
Cauchy--Schwarz. If
\begin{equation}\label{eq:cross}
 \E_{h\in[L]}
 \left|\E_{n\in[L]}F(n)G(n+h)\tau_h(n)\right|\ge\eta
\end{equation}
where \(F\) is 1-bounded and zero off \([L]\), \(G\) is
1-bounded on \([2L]\), and each \(\tau_h\) is a 1-bounded
degree-\(u\) nilsequence with dimension at most \(d\),
complexity at most \(M\), and Lipschitz norm at most \(M\),
then, when \(L\gg\eta^{-2}\), at least
\(c\eta^2L\) positive shifts \(a\) have supplied witnesses
\(\sigma_a\) with
\begin{equation}\label{eq:derivative}
 \left|\E_{n\in[L]}
       F(n)\overline{F(n+a)}\sigma_a(n)\right|\ge c\eta^2.
\end{equation}
Here \(\sigma_a\) is the product of two shifted
\(\tau_h\)'s, one conjugated, so its degree is unchanged
and its dimension at most doubles; its logarithmic
complexity and Lipschitz costs remain controlled.

For completeness, choose unit scalars \(\omega_h\)
aligning \eqref{eq:cross} and set
\[
 U(m)=\sum_{h=1}^L F(m-h)\omega_h\tau_h(m-h).
\]
Cauchy--Schwarz gives
\(\sum_m|U(m)|^2\ge\eta^2L^3/2\).
The diagonal terms contribute at most \(L^2\).
Pairing positive and negative differences in the
remaining sum yields
\[
 \sum_{a=1}^{L-1}\max_{1\le b\le L-a}
 \left|\frac1L\sum_n
 F(n)\overline{F(n+a)}
 \tau_{b+a}(n)\overline{\tau_b(n+a)}\right|
 \ge c'\eta^2L.
\]
Each summand is at most one, so a proportion
\(\gg\eta^2\) exceeds \(c\eta^2\). Choose \(b\)
separately for each \(a\); there is no loss for the
number of available witnesses. Unbounded shifts enter
the polynomial orbit, not the nilmanifold function.
If \(L\ll\eta^{-2}\), Parseval and
\(\E|F|^2\ge\eta^2\) instead give a linear-phase
correlation of size at least \(\eta/\sqrt L\).

The non-inverse construction of \cite[Lemma 6.2]{lss}
accepts the supplied family directly: omit the paragraph
that obtains the witnesses from the derivative Gowers
norm, and retain its normalization, structure-constant
pigeonhole, common-function selection, and vertical
Fourier steps. These use only \(d,M,\rho\).
The subsequent correlation-structure induction
(Lemma 6.3 through Theorem 6.4) and the integration
construction in Section 12 use no further ordinary
inverse theorem before its final Gowers/inverse steps.
Their budget is \(K_sx^{F_s}\).

Immediately before those final steps, the argument
supplies \eqref{eq:cross}, with
\[
 F(n)=f(sn+q_0)A(n),\qquad \deg A\le s,\qquad
 \deg\tau_h\le s-2.
\]
The known factor \(A\) includes the nilsequence
multiplier already present in that correlation.
Normalize all factors within the existing logarithmic
budget. Applying \eqref{eq:derivative} gives exactly
the supplied-family hypothesis for \(J_{s-1}\),
with input size at most \(K_sx^{F_s}\). Its output
\(\Theta\) gives the final test \(A\Theta\), of
degree at most \(s\), against \(f\) on the selected
residue class. The final residue-class Fourier
expansion in \cite{lss} returns the correlation on
\([N]\). In particular, the known factor \(A\)
is retained throughout.

Degree one is elementary. The base \(J_2\) follows from the fixed \(U^3\)
case after applying the converse bound to the supplied
linear witnesses. If a size condition fails at any
stage, \(\log N\le K_sx^{F_s}\); the available
correlation gives a lower bound on \(\E|f|^2\), and
Fourier analysis supplies the same logarithmic-budget
alternative.

Enlarge \(F_s\) within the class \eqref{eq:Fs} to
include initialization and final assembly. Sufficient
recurrences are
\begin{equation}\label{eq:IJ}
 J_s\le F_s(J_{s-1}+1),\qquad
 I_s\le F_s(J_s+1)(I_{s-1}+1).
\end{equation}
The second recurrence first uses the ordinary
degree-\((s-1)\) inverse theorem once to supply
derivative witnesses, then uses \(J_s\).
Taking logarithms and summing, with
\(\log F_s\ll s^c\) for an absolute \(c\), gives
\[
 \log(J_s+1)\ll s^{c+1},\qquad
 \log(I_s+1)\ll s^{c+2}.
\]
This proves the proposition.
\end{proof}

\subsection{The numerical library}

\begin{proposition}\label{prop:local}
For the construction of Sections~\ref{sec:macro}--\ref{sec:active},
each round introduces at most \(K_kp^F\) fresh absolute
coordinates, and the width increments, root logarithmic
costs, and logarithmic side requirements have a common
polynomial degree \(B\), with
\begin{equation}\label{eq:FB}
 1\le F\le B,\qquad 2\le B\le\exp(D^{O(1)}).
\end{equation}
The width increment at weight \(i\) uses only
width logarithms \(Q_j\) with \(j>i\).
\end{proposition}

\begin{proof}
Let \(x\ge2\) dominate the input dimensions and the
logarithms of rational heights, function complexities,
reciprocal errors, and scale requests. Simultaneous
rational elimination, bounded-step BCH, coefficient
arrays, Fourier truncation, and finite-dimensional
meshes have degree at most
\[
 H_D\le\exp(D^{O(1)})
\]
in \(x\). An external inverse call has degree at most
\(H_D(I+1)\), where
\(I=\max_{r\le D+2}I_r\).
Proposition~\ref{prop:inverse} gives
\(\log(I+1)\le D^{O(1)}\).

Two distinctions prevent artificial exponentiation.
First, a fixed auxiliary array, such as a Boolean
jet with at most \(2^{D+2}\) entries, contributes a
constant depending on \(D\). Its sublevel exponent
multiplies a logarithmic accuracy; it does not
iterate a budget map once per entry. A degree-\(d\)
array in a genuinely growing number \(n\) of
variables, by contrast, contributes the power
\((n+1)^d\), which is retained.

Second, the potentially long tuple loops in the
rank-transfer proof propagate sharpened bounds
\[
 d_{\rm out}\mid d_{\rm in}D_0^{c_D},\qquad
 B_{\rm out}\le B_{\rm in}\exp(P_D(n)).
\]
Their logarithms are additive in the previous
denominator and norm budgets. The cold costs of
proving a sharpened bound are not inputs to the
next sharpened output. Independent tensor children
start with the same parent budget and are combined
additively; new short thresholds descend in degree.
Products of probabilities or moduli likewise add
logarithmic costs.

The dependency chains can therefore be grouped as
follows. All references in the first column are to
\(\M\).

\begin{center}
\small
\renewcommand{\arraystretch}{1.14}
\begin{tabularx}{\linewidth}{@{}>{\raggedright\arraybackslash}p{0.29\linewidth}X@{}}
\toprule
Operations & Quantitative depth \\ \midrule
Rational calculus and conversion; Appendices A--B
 & Joint linear systems and at most \(D\) layer passages;
   one lower-degree child in recursive returns.\\
Comparison, counting, and residue transfer; Appendices C--E
 & Degree or axis passages; fixed-degree peeling and
   adaptive removals paid by telescoping potentials.\\
Samplers and globalization; Appendices F--I
 & Layer, grade, and old-weight passages; closed-form
   sublevel bounds and fixed sampler schedules.\\
Rank transfer and active return; Sections 4--8
 & Additive tuple/cut costs, simultaneous pivots, and
   one lower-degree child per return.\\
Extraction, compression, and persistent history;
Sections 9--10 and Sections~\ref{sec:macro}--\ref{sec:active} here
 & Macro, charge, and weight passages; joint coefficient
   arrays and finite meshes; descending numerical elimination.\\
\bottomrule
\end{tabularx}
\end{center}

Thus a fixed number of operation families have at most
\(D^{O(1)}\) successive nonlinear substitutions in
their combined dependency chains. The persistent
record adds no inverse invocation. Its two outer
factors share one child; integer evaluation menus
are meshed simultaneously; unitriangular blocks
are inverted by one matrix operation. There is
no iteration over the number of short values or
mesh members.

For the degree-one comparison, the explicit
bounded-potential argument of Appendix C has
polynomial logarithmic cost; the number of
peeling operations does not count as nonlinear
depth. Summing errors over axes, taking products
of selected denominators, and counting independent
alternatives likewise preserve polynomial
logarithmic budgets. Hence, after the structural
path-count substitutions and elimination of all
intermediate numerical records,
\[
 F,B\le[H_D(I+1)]^{D^{O(1)}}
      \le\exp(D^{O(1)}).
\]
The canonical-field and history construction in
Proposition~\ref{prop:macro} uses creation-layer
data only at strictly lower layers. The descending
schedule of Section 10.3 of \(\M\) therefore
retains precisely the stated width dependence.
\end{proof}

\section{The density-increment calculation}

We finish with the original dimension ledger and one
width recurrence.

\begin{proof}[Proof of Theorem~\ref{thm:main}]
Let \(F,B\) be as in Proposition~\ref{prop:local}.
The original rank-cut recurrence has sufficient
coordinate exponents
\[
 a_D=F+1,\qquad a_i=2a_{i+1}+1,
 \qquad
 a_i=(F+2)2^{D-i}-1.
\]
Allowing an extra factor \(p\) for cumulative cut
counts, a common structural exponent is
\[
 a=2^D(F+3).
\]
Thus all coordinate and cut-count inputs are
\(O_k(p^a)\).

Let \(v_i\) be an exponent for the width logarithm
at weight \(i\) after all \(T\ll_k p\) rounds.
The strict higher-weight dependence allows
\[
 v_D=1+Ba,\qquad v_i=1+Bv_{i+1}\quad(i<D).
\]
Each added \(1\) pays the sum over rounds.
Since \(B\ge2\) and \(a\ge1\),
\[
 v_1\le(2B)^Da.
\]
One final polynomial evaluation and summation
give total logarithmic cost and maximum
logarithmic side requirement \(O_k(p^A)\), where
\[
 A=1+Bv_1\le(2B)^{D+1}a.
\]
In particular,
\begin{equation}\label{eq:logA}
 \log A
 \le(D+1)\log(2B)+D\log2+\log(F+3)
 \le D^{O(1)}=k^{O(1)}.
\end{equation}

The fixed multiplicative density gain in \(\M\)
then gives a contradiction whenever
\[
 \log N\ge K_k\bigl(2+\log(1/\alpha)\bigr)^A.
\]
Inverting this condition yields
\[
 r_k(N)\le K'_kN
       \exp\!\left[-c'_k(\log N)^{1/A}\right].
\]
By \eqref{eq:logA}, enlarge one absolute constant
\(C\) so that \(A\le\exp(k^C)\) for every \(k\ge3\).
Decreasing the exponent to
\(\varepsilon_k=\exp(-k^C)\) and adjusting the
\(k\)-dependent constants gives the stated bound.
\end{proof}

\begin{thebibliography}{8}
\small
\bibitem{openai}
OpenAI, \emph{Quasipolynomial Bounds for Arithmetic
Progressions}, September 23, 2026, 198 pp.,
OpenAI/math, family 159.
\url{https://github.com/openai/math}.

\bibitem{lss}
J. Leng, A. Sah, and M. Sawhney,
\emph{Quasipolynomial bounds on the inverse theorem
for the Gowers \(U^{s+1}[N]\)-norm},
\href{https://arxiv.org/abs/2402.17994v3}{arXiv:2402.17994v3}.

\bibitem{leng}
J. Leng, \emph{Efficient Equidistribution of Nilsequences},
\href{https://arxiv.org/abs/2312.10772v5}{arXiv:2312.10772v5}.

\bibitem{lengquadratic}
J. Leng, \emph{Improved Quadratic Gowers Uniformity
for the M\"obius Function},
\href{https://arxiv.org/abs/2212.09635v3}{arXiv:2212.09635v3},
Theorem 8 and Appendix A.

\bibitem{jamneshan}
A. Jamneshan and T. Tao,
\emph{The inverse theorem for the \(U^3\) Gowers
uniformity norm on arbitrary finite abelian groups:
Fourier-analytic and ergodic approaches},
\href{https://arxiv.org/abs/2112.13759v3}{arXiv:2112.13759v3};
Discrete Analysis 2023:11.

\bibitem{milicevic}
L. Mili\'cevi\'c, \emph{A bilinear Bogolyubov argument
in abelian groups},
\href{https://arxiv.org/abs/2109.03093v3}{arXiv:2109.03093v3};
Discrete Analysis 2024:20, Proposition 28.

\bibitem{sanders}
T. Sanders, \emph{On the Bogolyubov--Ruzsa lemma},
Analysis \& PDE \textbf{5} (2012), no.~3, 627--655.
\end{thebibliography}
\end{document}
